\begin{entry}{Q017}{The AKLT state}

\entrymeta{
  topic   = {Matrix product states; AKLT model; valence bonds; parent Hamiltonians;
    addition of angular momentum.},
  est     = $\sim$50 minutes --- a multi-session entry. Natural split: (a)--(f), the
    construction; (g)--(i), the MPS and its parent Hamiltonian; (j) optional.,
  origin  = {Sourced. Questions~1--9 and the bonus of the source problem, restated with
    the basis ordering, the construction of the state and the Hamiltonian made explicit
    so the statement stands alone.},
  source  = {J.~E.~Nava~Sejas and D.~Rosa, \emph{The AKLT state}, lecture notes
    (minicourse), 21~September~2026. Original model:
    I.~Affleck, T.~Kennedy, E.~H.~Lieb and H.~Tasaki, Phys.~Rev.~Lett.~\textbf{59}, 799
    (1987).},
  prereq  = {Addition of angular momentum; the definition of an MPS.},
  status  = open,
}

\begin{problem}
\textbf{Setting.} Consider a chain of $N\ge 2$ spin-$1$ particles with periodic boundary
conditions, $N+1\equiv 1$. The local Hilbert space at site $j$ is $\CC^{3}$, with basis
$\{\ket{+_j},\ket{0_j},\ket{-_j}\}$. Each site is represented by a pair of
\emph{virtual} spin-$\tfrac12$ particles $(L_j,R_j)$, the ``left'' and ``right''
components, each with basis $\{\ket{\uparrow},\ket{\downarrow}\}$, ordered as
$(\uparrow,\downarrow)$ whenever an index is written as a matrix row or column. Write
$\sigma^{\pm}:=\tfrac12(\sigma^{x}\pm i\sigma^{y})$, and let $\hat S_j^{\alpha}$,
$\alpha=x,y,z$, be the spin-$1$ operators at site $j$.

\begin{enumerate}[label=\textbf{(\alph*)},leftmargin=*]
  \item Is the map from the virtual pair $(L_j,R_j)$ to the spin-$1$ site one-to-one?

  \item Given (a), adopt the identification
        \[
          \ket{+_j}\equiv\ket{\uparrow_j\uparrow_j},
          \qquad
          \ket{-_j}\equiv\ket{\downarrow_j\downarrow_j},
          \qquad
          \ket{0_j}\equiv\tfrac{1}{\sqrt2}\bigl(\ket{\uparrow_j\downarrow_j}
            +\ket{\downarrow_j\uparrow_j}\bigr).
        \]
        What is the physical motivation for this choice?

  \item Write the identification as a single operator
        $\hat P_j: \CC^{2}\otimes\CC^{2} \to \CC^{3}$ from the virtual space of site $j$ to
        its physical space. Verify that $\hat P_j\hat P_j^{\dagger}=\id_{3}$, that
        $\hat P_j^{\dagger}\hat P_j$ is the orthogonal projector onto the triplet
        subspace of $L_j\otimes R_j$, and identify $\ker\hat P_j$.

  \item Define the matrix elements
        $P^{(j)s}_{ab}:=\bra{s_j}\hat P_j\ket{a_j b_j}$, with
        $s_j\in\{+_j,0_j,-_j\}$ and $a_j,b_j\in\{\uparrow,\downarrow\}$. Show that
        \[
          P^{(j)+}_{ab}=\delta_{a\uparrow}\delta_{b\uparrow},
          \qquad
          P^{(j)0}_{ab}=\tfrac{1}{\sqrt2}\bigl(\delta_{a\uparrow}\delta_{b\downarrow}
            +\delta_{a\downarrow}\delta_{b\uparrow}\bigr),
          \qquad
          P^{(j)-}_{ab}=\delta_{a\downarrow}\delta_{b\downarrow}.
        \]

  \item The right virtual spin of site $j$ and the left virtual spin of site $j+1$ are
        required to form the singlet
        \[
          \ket{\omega}_{R_j,L_{j+1}}
          =\tfrac{1}{\sqrt2}\bigl(\ket{\uparrow}_{R_j}\ket{\downarrow}_{L_{j+1}}
            -\ket{\downarrow}_{R_j}\ket{\uparrow}_{L_{j+1}}\bigr).
        \]
        Show that this can be written as
        \[
          \ket{\omega}_{R_j,L_{j+1}}
          =\tfrac{1}{\sqrt2}\sum_{b,c\in\{\uparrow,\downarrow\}}\epsilon_{bc}\,
            \ket{b}_{R_j}\ket{c}_{L_{j+1}},
          \qquad
          \epsilon=i\sigma^{y},
        \]
        and check that $\epsilon$ is the fully antisymmetric tensor in two dimensions.

  \item The state of the chain is built by first placing every bond $(R_j,L_{j+1})$ in
        the singlet, then mapping each site to its physical space:
        \[
          \ket{\Phi}
          :=\Bigl(\textstyle\bigotimes_{j=1}^{N}\hat P_j\Bigr)
            \Bigl(\textstyle\bigotimes_{j=1}^{N}\ket{\omega}_{R_j,L_{j+1}}\Bigr).
        \]
        Argue that the two constraints --- the on-site map $\hat P_j$ and the bond
        singlet --- combine, site by site, into the single tensor
        \[
          A^{[j]s}_{ac}=\tfrac{1}{\sqrt2}\sum_{b\in\{\uparrow,\downarrow\}}
            P^{(j)s}_{ab}\,\epsilon_{bc},
        \]
        and state which virtual spin each of the indices $a$ and $c$ refers to.

  \item Show that
        \[
          \braket{s_1s_2\cdots s_N}{\Phi}
          =\Tr\bigl[A^{s_1}A^{s_2}\cdots A^{s_N}\bigr],
        \]
        where $A^{s}$ denotes the $2\times2$ matrix $(A^{[j]s}_{ac})_{a,c}$, which does
        not depend on $j$. Conclude that $\ket{\Phi}$ is a translation-invariant MPS of
        bond dimension $\chi=2$, and derive
        \[
          A^{+}=\tfrac{1}{\sqrt2}\,\sigma^{+},
          \qquad
          A^{0}=-\tfrac12\,\sigma^{z},
          \qquad
          A^{-}=-\tfrac{1}{\sqrt2}\,\sigma^{-}.
        \]

  \item Consider the Hamiltonian
        \[
          \hat H=\sum_{j=1}^{N}\hat\Pi^{S=2}_{j,j+1},
        \]
        where $\hat\Pi^{S=2}_{j,j+1}$ is the hermitian projector onto the subspace of
        $\CC^{3}_j\otimes\CC^{3}_{j+1}$ on which
        $(\hat{\vec S}_j+\hat{\vec S}_{j+1})^{2}=2(2+1)$, i.e.\ total spin $2$, tensored
        with the identity on all other sites. Argue that the spectrum of $\hat H$ is
        non-negative.

  \item Using the construction in terms of the virtual spins, show that
        \[
          \hat\Pi^{S=2}_{j,j+1}\ket{\Phi}=0\quad\text{for all } j,
          \qquad\text{hence}\qquad
          \hat H\ket{\Phi}=0 ,
        \]
        and conclude, using (h), that $\ket{\Phi}$ is a ground state of $\hat H$. (You may take for
        granted that $\ket{\Phi}\neq 0$ for $N\ge 2$.)

  \item \emph{Bonus.} Show that $\hat H$ can be rewritten as
        \[
          \hat H=\frac16\sum_{i=1}^{N}\Bigl[\bigl(\hat{\vec S}_i\cdot\hat{\vec S}_{i+1}
            \bigr)^{2}+3\,\hat{\vec S}_i\cdot\hat{\vec S}_{i+1}+2\Bigr].
        \]
\end{enumerate}
\end{problem}

% ============================================================== solutions ===

\begin{solution}[Part (a)]
The two spin-1/2 particles live in a 4-dimensional space, meanwhile, for one spin-1 live in a 3-dimensional one, hence the map is not injective
\end{solution}

\begin{solution}[Part (b)]
It follows usual Clebsch-Gordan coefficients for spin summation from different particle subspaces.
\end{solution}

\begin{solution}[Part (c)]
\begin{equation}
    P_j\ket{\uparrow_L\uparrow_R}_j = \ket{+_j} = \begin{pmatrix}
        1\\
        0\\
        0
    \end{pmatrix}\quad P_j \ket{\downarrow_L\downarrow_R}_j = \ket{-_j} = \begin{pmatrix}
        0\\
        0\\
        1
    \end{pmatrix} \quad P_j\frac{1}{\sqrt2}\bigl(\ket{\uparrow_j\downarrow_j}
            +\ket{\downarrow_j\uparrow_j}\bigr). = \ket{0_j} = \begin{pmatrix}
        0\\
        1\\
        0
    \end{pmatrix}
\end{equation}
\end{solution}

\begin{solution}[Part (d)]
We can see that
\begin{equation}
    P_j = \begin{pmatrix}
    \bra{\uparrow_L\uparrow_R}_j\\
    \frac{1}{\sqrt{2}}(\bra{\uparrow_L\downarrow_R}_j+\bra{\downarrow_L\uparrow_R}_j)\\
    \bra{\downarrow_L\downarrow_R}_j
    \end{pmatrix}= \begin{pmatrix}
    1 & 0 & 0 &0\\
    0 & \frac{1}{\sqrt{2}} & \frac{1}{\sqrt{2}} &0\\
    0 & 0 & 0 &1
    \end{pmatrix}
\end{equation}


Just take the hermitian conjugate to see
\begin{equation}
    P_j^\dagger= \begin{pmatrix}
    1 & 0 & 0 &\\
    0 & \frac{1}{\sqrt{2}} & 0 \\
    0 & \frac{1}{\sqrt{2}} & 0 & \\
    0 & 0  & 1
    \end{pmatrix}
\end{equation}
that their multiplication leads to the identity.

\end{solution}

\begin{solution}[Part (e)]
Follows directly from the summation expression
\end{solution}

\begin{solution}[Part (f)]
We can think constructively, where $P_j$ is an action on each site and anti-symmetric tensor $\epsilon_{a,b}$ determines the link between different sites. In that sense we can think of a tensor connecting $L_j$ to $L_{j+i}$ that represents the projection and the link between sites, naturally represented by $P^{(j)s}$ contracted with $\epsilon$ on the degrees of freedom corresponding to $R_j$
\end{solution}

\begin{solution}[Part (g)]
Let
$$\ket{(a,b)}=\ket{a_1b_1a_2b_2...a_nb_n}$$ for $a,b \in \{+,-\}^N$, such that $a_i$ refers to the $L_i$ spin-1/2 virtual particle and $b_i$ refers to $R_i$ spin-1/2 virtual particle. Also take $s\in \{-1,0,1\}^N$
Using that notation, the matrix elements of the tensor product of the projector can be written as 
$$\avg{s}{\bigotimes_{j=1}^N P_j}{(a,b)}= \prod_j P^{(j)s}_{a_j,b_j}$$
whereas, considering the periodic boundary condition where N+1=1,
$$\bra{(a,b)}\bigotimes_{j=1}^N\ket{\omega_j}= \frac{1}{\sqrt{2^N}}\prod_i \epsilon_{b_i,a_{i+1}}$$
The state $\ket{\Phi}$ lives in the space of N spin-1 particles, hence can be written as 
\begin{equation}
    \ket{\Phi} = \sum_{s } c_s \ket{s}
\end{equation}
where the coefficients are given by 
\begin{equation}
\begin{split}
    c_s =& \bra{s}\ket{\Phi}\\
    =& \bra{s} \left( \bigotimes_{j=1}^N P_j \right) \sum_{(a,b)} \proj{(a,b)} \left(\bigotimes_{l=1}^N\ket{\omega_l}\right)\\
    =& \sum_{a,b \in \{+,-\}^N}\prod_jP^{(j)s}_{a_j,b_j} \prod_l \epsilon_{a_l,b_{l+1}}\\
    =&(-1)^N \sum_{a,b \in \{+,-\}^N}\prod_jP^{(j)s}_{a_j,b_j} \epsilon_{b_j,a_{j-1}}\\
    =& (-1)^N \sum_{a} \prod_j (A^{[j]s})_{a_j,a_{j-1}}\\
    =& \Tr \left( \prod_j(A^{[j]s})\right)
\end{split}
\end{equation}
Now, we see that each matrix $A^{[j]s}_{a,c}$ is a 2x2 matrix establishing the bound dimension. Moreover, the cyclic property of the trace provides the translation invariant property on spin-1 lattice sites on the chain. 

By its expression

\begin{equation}
    A^{[j]+}_{ac}=\tfrac{1}{\sqrt2}\sum_{b\in\{\uparrow,\downarrow\}}
            \delta_{a\uparrow}\delta_{b \downarrow}\,\epsilon_{bc} = \frac{1}{\sqrt{2}}\delta_{a\uparrow}\epsilon_{\downarrow c}
\end{equation}
For the row $a= \downarrow$ the right-hand side is always zero. For $a=\uparrow$, the only entry that survives is when $c=\uparrow$, so
\begin{equation}
A^{+} = \frac{1}{\sqrt{2}} \sigma^+
\end{equation}

The other cases follow similarly.



\end{solution}




\begin{solution}[Part (h)]
Each Hermitian projector is a positive operator, hence their sum (the Hamiltonian) is positive also.
\end{solution}

\begin{solution}[Part (i)]

The spin-2 sector generated from the sum of spin-1 particles on sites $j$ and $j+1$, can be rewritten as the sum of four virtual spin-1/2 particles $L_j,R_j,L_{j+1},R_{j+1}$. The projector $\Pi^{S=2}_{j,j+1}$ is an operator on the two spin-1 particles, the corresponding projector on the four virtual spin-1/2 particle is denoted $$\Pi^{S=2}_{L_j,R_j,L_{j+1},R_{j+1}} \equiv \Pi^{S=2}_{4\ \text{spin-1/2}},$$ such that, being $\hat P_j: \CC^{2}\otimes\CC^{2} \to \CC^{3}$ just an identification, 
\begin{equation}
   (P_jP_{j+1}) \circ \Pi^{S=2}_{L_j,R_j,L_{j+1},R_{j+1}}  =\Pi^{S=2}_{j,j+1} \circ(P_jP_{j+1}) : \CC^{2}\otimes\CC^{2} \to \CC^{3}  
\end{equation}

Consider 
\begin{equation} 
    \ket{\Omega}= \bigotimes_{l=1}^N\ket{\omega_l}.
\end{equation}
Verifying that $\hat\Pi^{S=2}_{j,j+1}\ket{\Phi}=0$ is equivalent to $\Pi^{S=2}_{4\ \text{spin-1/2}}\ket{\Omega}=0$, by the previous identification

To proceed, one must notice that the singlet has total spin-0. So the only contribution for the total spin for the four virtual spin-1/2 particles $L_j,R_j,L_{j+1},R_{j+1}$, is from $L_j$ and $R_{j+1}$ and the sum of two spin-1/2 particles can generate at most a total spin 1, thus, the total sum of spin from all four virtual particles can never achieve $2$. Therefore, $\hat\Pi^{S=2}_{j,j+1}\ket{\Phi}=\Pi^{S=2}_{4\ \text{spin-1/2}}\ket{\Omega}=0$ as desired.

\end{solution}

\begin{solution}[Part (j)]
Let $\vec{J}_i= \vec{S}_i + \vec{S}_{i+1}$ be the total spin of the sum of spin-1 particles from two neighbours sites. Now, from the expression of the total spin
\begin{equation}
\begin{split}
    J_i^2 &= ( \vec{S}_i + \vec{S}_{i+1})\cdot( \vec{S}_i + \vec{S}_{i+1})\\
    &= \vec{S}_i^2 + 2  \vec{S}_i\cdot \vec{S}_{i+1}  + \vec{S}_{i+1}^2\\
    &= 4+ 2  \vec{S}_i\cdot \vec{S}_{i+1}
\end{split}
\end{equation}
where in the last line it was used that for spin-1 particles $\vec{S}_{i+1}^2=1(1+1)=2$. 

The eigenvalues of the total spin $J_i^2$ are $j_i(j_i+1)$ for $j_i=0,1,2$. Now, for an eigenstate with $j_i=0$, the total spin itself already nullifies that component, leaving only the direction associated with $j_i=1,2$. Nonetheless, in general you can apply the same reasoning to notice that $J_i^2-j_i(j_i+1)\id$ is an operator that nullifies the eigenspace associated with the eigenvalue $j_i(j_i+1)$ of $J_i$.

So in order to construct an projector orthogonal both to the total spin 1 and total spin 0, one can just consider the operator $J_i^2 (J_i^2 - 2)$ ($J_i$ forces orthogonality with spin-0 and $J_i^2 - 2$ to spin-1), yet, one must normalize for $j_i=2$ to construct an actual total spin-2 projector

Therefore,
\begin{equation}
    \begin{split}
        \Pi^{S=2}_{i,i+1} &= \frac{1}{6\cdot4} J_i^2 (J_i^2 - 2)\\
        &= \frac{1}{6\cdot4} \left( 4+ 2  \vec{S}_i\cdot \vec{S}_{i+1}\right) \left(2+ 2  \vec{S}_i\cdot \vec{S}_{i+1}\right)\\
        &= \frac{1}{6}\left( 2 + 3\vec{S}_i\cdot \vec{S}_{i+1} + \left(\vec{S}_i\cdot \vec{S}_{i+1}\right)^2\right).
    \end{split}
\end{equation}
The expression for the Hamiltonian follows directly

\end{solution}



\attempt{
  date        = 2026-09-21,
  time        = took days to finish it intertwined with others and the school,
  difficulty  = 3,
  how it went = ,
}

\end{entry}